% ECONOMICS_PROBLEM: {"id": "C73-564", "title": "Learning-based equilibrium selection", "jel_code": "C73", "source_id": "OP-0281"}
% !TeX program = xelatex
\documentclass[11pt,a4paper]{article}
\usepackage[margin=23mm]{geometry}
\usepackage{fontspec}
\setmainfont{Latin Modern Roman}
\usepackage{amsmath,amssymb,mathtools}
\usepackage{xcolor,enumitem,booktabs,longtable,array,xurl,hyperref}
\definecolor{muted}{HTML}{53636C}
\hypersetup{colorlinks=true,urlcolor=blue,linkcolor=blue}
\setlength{\parindent}{0pt}
\setlength{\parskip}{5pt}
\newcommand{\R}{\mathbb R}
\newcommand{\E}{\mathbb E}
\newcommand{\N}{\mathbb N}
\newcommand{\ind}{\mathbf 1}
\newcommand{\Var}{\operatorname{Var}}
\newcommand{\Cov}{\operatorname{Cov}}
\newcommand{\supp}{\operatorname{supp}}
\DeclareMathOperator*{\argmin}{arg\,min}
\DeclareMathOperator*{\argmax}{arg\,max}
\newcommand{\entrymeta}[3]{{\sffamily\small\color{muted}\textbf{\texttt{#1}}\quad|\quad #2\par #3\par}}
\newcommand{\Problemref}[1]{\texttt{#1}}
\newcommand{\JELref}[2]{\texttt{#2} (JEL \texttt{#1})}
\begin{document}
\section*{Learning-based equilibrium selection}
\textbf{Problem ID:} \texttt{C73-564}\quad \textbf{JEL:} \texttt{C73}\par
% BEGIN ECONOMICS STATEMENT
\label{problem:OP-0281}
{\sffamily\small\color{muted}Original subject group: Game theory and strategic interaction\par}
\label{definitions:problem:30007563}
\textbf{Audit classification:} Published research agenda. Fudenberg–Levine pp.153–154 motivates equilibrium prediction through learning. The finite-horizon selection model is editorial.
\subsection*{Definitions and precise research target}
\textbf{Model and research target.} Consider a preregistered class of finite coordination games \(G\) with a known set \(E_G\) of Nash equilibria. Groups of human players repeatedly interact for \(T\) periods under randomly assigned feedback and initial-history treatment \(Z\). Let \(\widehat\sigma_T\) be the product of empirical action distributions, and fix disjoint neighborhoods \(U_e\) around isolated equilibria \(e\in E_G\). Define the selection probabilities
\[
 q_e(G,Z,P,T)=\Pr_P(\widehat\sigma_T\in U_e\mid do(Z)),
\]
where \(P\) is the sampling population; preserve a residual category for paths lying outside every neighborhood. The target is to estimate these probabilities and identify a portable model relating them to learning-rule heterogeneity, payoff risk, information, and initial experience. Separate a treatment's effect on equilibrium selection from its effect on the speed of approaching any equilibrium by measuring several prespecified horizons. Elicit individual beliefs or learning parameters with measurement error explicitly modeled, and test the selection rule on new games and populations. Random assignment identifies treatment effects within sampled groups; transporting results requires stated population-overlap assumptions.

\emph{Definitions and maintained assumptions (editorial unless a source definition is named).} For each coordination game select a finite set \(E_G\) of isolated Nash equilibria; a game with a continuum of equilibria requires a different target. Use the \(\ell^1\) metric on product strategy space and radii \(r_e>0\) making \(U_e=\{\sigma:\|\sigma-e\|_1<r_e\}\) pairwise disjoint. The event probabilities refer to independently sampled groups in population \(P\) assigned treatment \(Z\); \(do(Z)\) is this randomized group assignment. Define the residual probability \(q_0=1-\sum_e q_e\). It includes nonconvergence, other equilibria and points on the neighborhood boundaries. Specify empirical profile construction under fixed roles and the number of periods \(T\). A portable model is a prespecified map from game features, treatment and participant measurements to the entire vector \((q_0,(q_e)_e)\), evaluated by a common proper scoring rule on unseen games. Randomized treatment identifies sampled-population selection probabilities, not structural learning heterogeneity without a measurement model.
\subsection*{Variants and assumption changes}
\begin{enumerate}
\item \textbf{Two pure equilibria (editorial).} Fix two-by-two coordination games and vary payoff risk and initial play independently. Estimate changes in basin-selection frequencies and time to enter either neighborhood.
\item \textbf{Heterogeneous learners (editorial).} Admit a finite mixture of specified learning rules and imperfect belief reports. Determine which treatments identify mixture weights and which selection predictions remain partially identified when distinct rules generate the same history law.
\end{enumerate}
\subsection*{References and source audit}
\begin{enumerate}
\item pp.153–154 and conclusion pp.165–166. Supports the stated research area; exact displayed specification is editorial. \url{https://economics.mit.edu/sites/default/files/2022-09/Whither%20Game%20Theory%20Towards%20a%20Theory%20of%20Learning%20in%20Games.pdf}
\end{enumerate}
\begin{enumerate}\raggedright
\item\hypertarget{definitions:source:30007563:1}{} Drew Fudenberg; David K. Levine. 2016. \emph{Whither Game Theory? Towards a Theory of Learning in Games}. Introduction, p.153; conclusion, pp.165–166.
\url{https://economics.mit.edu/sites/default/files/2022-09/Whither%20Game%20Theory%20Towards%20a%20Theory%20of%20Learning%20in%20Games.pdf}.
DOI: \url{https://doi.org/10.1257/jep.30.4.151}.
\end{enumerate}

\subsection*{Common conventions: Source mathematical conventions}

These conventions apply to each entry unless explicitly replaced.

\textbf{Models and identification.} A primitive model \(m\) specifies all probability laws, feasible choices, information, equilibrium and measurement rules used by its target. The admissible class \(\mathcal M\) is fixed independently of outcomes. If \(P_O(m)\) is its observable probability law and \(\tau(m)\) the target, its identified set at observable law \(P\) is
\[
 \mathcal I_\tau(P)=\{\tau(m):m\in\mathcal M,\ P_O(m)=P\}.
\]
Point identification means that this set is a singleton. An empty set signals incompatibility of data and assumptions. Coordinate bounds need not describe the joint set. Bounds are sharp only if every retained target is attainable or arbitrarily approachable by compatible models. Finite bounds require explicit restrictions on unobserved quantities; unrestricted confounding, hidden wealth, extrapolation or arbitrary equilibrium selection can yield uninformative sets.

\textbf{Causal interventions.} A potential outcome \(Y_i(p)\) is the outcome under a complete policy regime \(p\), including timing, financing, information and market spillovers. The target population and units are fixed before treatment. With interference, use the complete assignment vector or a justified exposure mapping; individual no-interference notation alone cannot identify an economy-wide intervention. A difference of mean potential outcomes does not require the cross-world joint distribution. Conditioning on post-treatment migration, survival, enrollment or employment changes the estimand and needs separate justification.

\textbf{Statistical statements.} An estimator, confidence set, forecast or test specifies a sampling law, model class and loss/coverage criterion. Uniform \((1-\alpha)\)-coverage means \(\inf_{m\in\mathcal M}\Pr_m\{\tau(m)\in C_n\}\geq1-\alpha\), or an explicitly asymptotic version. The whole parameter space trivially has coverage, so informativeness requires a stated width, risk or power objective. A forecasting improvement is relative to a named benchmark, information set and evaluation population.

\textbf{Equilibrium and optimization.} An equilibrium comprises feasible plans, optimality, beliefs where needed, budget constraints and all market/resource clearing. The primitives or a stated benchmark define each correspondence \(\mathcal E(p;m)\). Nonemptiness is assumed or proved, never inferred just from compact policy sets. A selected-equilibrium target uses a declared selection rule; a robust target quantifies over all permitted equilibria. Use suprema or infima unless attainment is established. An \(\varepsilon\)-optimal policy is feasible and within \(\varepsilon\) of the finite optimum value. Report an empty feasible set separately.

\textbf{Welfare and accounting.} Normative utility, interpersonal normalization, social weights, numeraire, discounting and the use of public receipts are primitives. Behavioral utility need not coincide with the welfare criterion. Household welfare includes ownership income and rents once; adding profits again to compensating variation generally double-counts them. A monetary equivalent is defined by an explicit transfer and price convention, with monotonicity and range conditions. Average effects, mechanism contributions, transfers and welfare losses are different targets. Infinite sums require convergence and dynamic feasible plans require terminal or no-Ponzi conditions.

\textbf{Source versions.} A working-paper date, revision date and journal publication year are distinguished. A DOI for a journal version is not automatically the DOI of an earlier manuscript. Locators refer to the stated version and printed pages where specified. A metadata-only check does not verify a claim about a full-text conclusion. Every entry's source subsection narrows the provenance claim accordingly.
% END ECONOMICS STATEMENT
\end{document}
